\setcounter{equation}{0}
\chapter{Introduction}
\label{chap:introduction}

\section{Overview}
\label{sec:overview}

This project builds a weather forecast for the SJCET campus from two very different kinds of data: a handwritten weather register kept at the college for close to twenty years, and decades of freely available regional weather data. The register is first digitized into a verified daily dataset. A graph neural network then learns how weather evolves across India from the regional data, learns how real weather stations differ from that data, and is finally fine-tuned on the digitized register. This staged \textbf{transfer learning} lets a modern forecasting model work at a single station whose own record is far too short to train one from scratch. The result is a six-day forecast of maximum and minimum temperature and rainfall at SJCET, served through a web application.

\section{Background}
\label{sec:background}

Weather shapes daily life in Kerala more directly than in most places. The state receives two monsoons, and the Western Ghats produce sharp local differences in temperature and rainfall over a few tens of kilometres, so a forecast that is good for a district can be poor for a particular campus. Public forecasts, however, are issued only for districts and grid cells.

The idea for the project began with the college's own weather station. For close to twenty years, staff at SJCET have read its instruments every working morning at about 08:30 and written the dry-bulb and wet-bulb temperatures, the maximum and minimum temperatures and the rainfall into a paper register, together with monthly rainfall totals. This is exactly the local record a campus forecast needs, but on paper it cannot be analysed or used to train a model. Hawkins et al.\ show that rescuing even a modest number of such paper observations changes what is known about past weather \cite{ref1}, so digitizing the register became the first goal.

Digitizing the register, however, does not by itself make a forecast possible. The clean scans cover July 2010 to December 2016, which gives only 1,549 observed days, because the register is read on working days only. Data-driven models such as Pangu-Weather \cite{ref5} and FuXi \cite{ref6} now rival operational weather prediction, but they are trained on decades of dense, gridded global data, and a network trained on a few years of one station's readings would simply memorise them. Regional reanalysis and station archives, on the other hand, cover India daily for decades, but describe grid-cell averages or other places rather than the campus.

Transfer learning connects the two. It reuses what a model learns on a data-rich source task to improve a related, data-poor target task \cite{ref8}, and Nipen et al.\ have shown that a graph neural network pretrained on global reanalysis and refined on a few years of regional data outperforms an operational regional forecast \cite{ref7}. Taking this idea one step further down, from region to a single station, led to the project: learn the weather from abundant regional data, and use the small, carefully digitized local record only to adapt the model to SJCET.

\section{Proposed Approach}
\label{sec:approach}

The project is carried out in five steps, which also form its objectives:

\begin{enumerate}[itemsep=0pt,topsep=4pt]
    \item \textbf{Digitize the register.} Transcribe every page, check it automatically and reconcile each month against the observers' rainfall totals, giving a traceable daily dataset.
    \item \textbf{Build the regional dataset.} Combine 36 years of NASA POWER reanalysis over India and Kerala, six NOAA weather stations around Kerala and the SJCET register into a single weather graph.
    \item \textbf{Train in three stages.} Pretrain a graph neural network on the regional reanalysis, refine it on the reference stations, and finally fine-tune it on the SJCET register.
    \item \textbf{Evaluate accuracy and data efficiency.} Measure forecast accuracy at SJCET on held-out years against conventional methods, and how much local data the final stage needs.
    \item \textbf{Deliver the forecast.} Present the forecasts, the register and the results in a web application.
\end{enumerate}

\lfoot{\textit{Department of Artificial Intelligence and Data Science, SJCET Palai}}
\renewcommand{\footrulewidth}{0.4pt}
